\section{Endliche Automaten}

\subsection{Grundidee}
Endliche Automaten (EA) sind speicherbeschränkte Berechnungsmodelle zur Erkennung formaler Sprachen.
Der einzige interne Speicher ist der aktuelle Zustand.
Ein Eingabewort wird sequenziell gelesen; Akzeptanz wird über Endzusände entschieden.

\subsection{Deterministische endliche Automaten}
Ein deterministischer endlicher Automat (DEA) ist ein Tupel
\[
    M = (Q, \Sigma, \delta, q_0, F)
\]
mit
\begin{itemize}
    \item endlicher Zustandsmenge \(Q\),
    \item Alphabet \(\Sigma\),
    \item totaler Übergangsfunktion \(\delta : Q \times \Sigma \to Q\),
    \item Startzustand \(q_0 \in Q\),
    \item Endzustandsmenge \(F \subseteq Q\).
\end{itemize}

\subsection{Konfigurationen}
Eine Konfiguration ist ein Paar
\[
    (q,w) \in Q \times \Sigma^*
\]
mit aktuellem Zustand \(q\) und Restwort \(w\).

\subsection{Berechnungsschritte}
Die Ein-Schritt-Relation ist definiert durch
\[
    (q,w) \vdash_M (p,x)
\]
genau dann, wenn \(w=ax\) mit \(a\in\Sigma\) und \(\delta(q,a)=p\).

Die reflexiv-transitive Hülle wird mit
\[
    (q,w) \vdash_M^* (p,x)
\]
bezeichnet.

\subsection{Akzeptanz und Sprache}
Für \(w\in\Sigma^*\) startet die Berechnung in \((q_0,w)\).
\(w\) wird akzeptiert, falls
\[
    (q_0,w) \vdash_M^* (q_f,\varepsilon) \quad\text{für ein } q_f\in F.
\]

Die erkannte Sprache ist
\[
    L(M)=\{\,w\in\Sigma^* \mid (q_0,w)\vdash_M^*(q_f,\varepsilon) \text{ für ein } q_f\in F\,\}.
\]

\subsection{Nichtdeterministische endliche Automaten}
Ein nichtdeterministischer endlicher Automat (NEA) hat dieselbe Struktur
\[
    M = (Q,\Sigma,\delta,q_0,F)
\]
wobei
\[
    \delta : Q \times \Sigma \to \mathcal{P}(Q)
\]
gilt. Damit sind mehrere Folgezustände pro Symbol möglich.

Akzeptanz: Es existiert mindestens ein Lauf, der in einem Zustand aus \(F\) endet.

\subsection{Äquivalenz von DEA und NEA}
Sein NEA \(N=(Q_N,\Sigma,\delta_N,q_{0N},F_N)\) kann in einen equivalenten DEA \(D=(Q_D,\Sigma,\delta_D,q_{0D},F_D)\) überführt werden durch
\begin{itemize}
    \item \(Q_D = \mathcal{P}(Q_N)\) (Mengen von NEA-Zuständen),
    \item \(F_D = \{S \subseteq Q_N \mid S \cap F_N \neq \emptyset\}\) (DEA-Endzustände sind Mengen, die mindestens einen NEA-Endzustand enthalten),
    \item \(\delta_D(S,a) = \bigcup_{q \in S} \delta_N(q,a)\) für \(S \subseteq Q_N\) und \(a \in \Sigma\) (Menge aller Zustände, die durch \(a\) von einem Zustand in \(S\) erreichbar sind),
\end{itemize}

\subsection{$\varepsilon$-Übergänge}
\(\varepsilon\)-NEAs erlauben Übergänge ohne Symbolkonsum mit
\[
    \delta : Q \times (\Sigma \cup \{\varepsilon\}) \to \mathcal{P}(Q).
\]
Auch \(\varepsilon\)-NEAs sind durch Eliminierung von \(\varepsilon\)-Übergängen und anschliessende Determinisierung zu DEA äquivalent.


\subsection{Äquivalenz DEA und RA}
Jede von einem DEA erkannte Sprache kann durch einen regulären Ausdruck beschrieben werden, und umgekehrt. Diese Sprachen werden als reguläre bezeichnet.
Reguläre Sprachen sind durch verschiedene Formalismen äquivalent beschreibbar, darunter DEA, NEA, \(\varepsilon\)-NEA und reguläre Ausdrücke.
\begin{itemize}
    \item \textbf{Akzeptierende Mech.:} DEA, NEA, \(\varepsilon\)-NEA
    \item \textbf{Beschreibungsform:} reguläre Ausdrücke
\end{itemize}

\subsection*{Abschlusseigenschaften regulärer Sprachen}
Die Klasse der regulären Sprachen ist unter verschiedenen Operationen abgeschlossen.

Seien $L, L_1, L_2 \subseteq \Sigma^*$ reguläre Sprachen. Dann sind ebenfalls regulär:
\begin{itemize}
    \item Vereinigung: \(L_1 \cup L_2\)
    \item Konkatenation: \(L_1 L_2\)
    \item Kleenesche Hülle: \(L^*\)
    \item Komplement: \(\overline{L} = \Sigma^* \setminus L\)
    \item Schnitt: \(L_1 \cap L_2\)
    \item Mengendifferenz: \(L_1 \setminus L_2\)
\end{itemize}

\subsection*{Zustandsklassen}
Zu jedem Zustand $p$ eines Automaten $M$ gehört die Zustandsklasse
\[
[p] = \{ w \in \Sigma^* \mid M \text{ endet nach dem Lesen von } w \text{ in } p \}.
\]

Die Zustandsklassen besitzen folgende Eigenschaften:
\begin{itemize}
    \item Die Menge aller Wörter wird durch die Klassen partitioniert:
    \[
    \Sigma^* = \bigcup_{p \in Q} [p]
    \]
    \item Für verschiedene Zustände sind die Klassen disjunkt:
    \[
    [p] \cap [q] = \emptyset \quad \text{für } p \ne q
    \]
    \item Die von $M$ akzeptierte Sprache ergibt sich als Vereinigung der Klassen akzeptierender Zustände:
    \[
    L(M) = \bigcup_{p \in F} [p]
    \]
\end{itemize}

\subsection*{Grenzen endlicher Automaten}
Ein endlicher Automat hat nur endlich viele Zustände und deshalb nur begrenzten Speicher.

Landen zwei Wörter $x$ und $y$ im selben Zustand, kann der Automat sie nicht mehr unterscheiden. Daher gilt für jedes Fortsetzungswort $z \in \Sigma^*$:
\[
xz \in L(M) \iff yz \in L(M).
\]

Wörter derselben Zustandsklasse verhalten sich also bei jeder Fortsetzung gleich.
Damit kann man Mindestzahlen für die benötigten Zustände eines Automaten ableiten.
